\subsection{Molekülorbitale und Elektronendichte}
\label{subsec:molekuelorbitale-elektronendichte}

Die bisherige Beschreibung chemischer Bindungen war bewusst anschaulich:
Atome geben Elektronen ab, nehmen Elektronen auf oder teilen sich
Elektronenpaare. Dieses Bild ist nützlich, aber es bleibt eine Vereinfachung.
In Molekülen sind Elektronen nicht einfach kleine Teilchen, die sich an einem
festen Ort zwischen zwei Atomen befinden. Ihre Verteilung muss genauer
beschrieben werden.

In der Quantenmechanik werden Elektronen durch Zustände beschrieben, die man
Orbitale\index{Orbital} nennt. Ein Orbital ist keine feste Bahn eines Elektrons
um den Atomkern. Es beschreibt vielmehr, in welchen Raumbereichen sich ein
Elektron mit hoher Wahrscheinlichkeit aufhält.

Bei einzelnen Atomen spricht man von Atomorbitalen\index{Atomorbital}. Nähern
sich zwei oder mehrere Atome einander an, so können sich ihre Atomorbitale
überlagern. Dadurch entstehen Molekülorbitale\index{Molekülorbital}. Diese
Molekülorbitale gehören nicht mehr nur zu einem einzelnen Atom, sondern zum
gesamten Molekül.

\begin{DidacticBox}[Ein Orbital ist keine Kreisbahn]
	\label{box:orbital-keine-kreisbahn}
	Ein Orbital darf nicht mit einer klassischen Bahn verwechselt werden. Ein
	Elektron bewegt sich im Molekül nicht wie ein Planet auf einer festen Bahn.
	
	Ein Orbital beschreibt eine räumliche Wahrscheinlichkeitsverteilung. Es zeigt,
	wo sich ein Elektron mit hoher Wahrscheinlichkeit aufhalten kann.
\end{DidacticBox}

Besonders wichtig ist die Elektronendichte\index{Elektronendichte}. Sie gibt
an, wie stark Elektronen in einem bestimmten Raumbereich vertreten sind. In
einer chemischen Bindung befindet sich häufig erhöhte Elektronendichte zwischen
den beteiligten Atomkernen. Diese Elektronendichte trägt zur Bindung bei, weil
sie die positiv geladenen Kerne gemeinsam anzieht.

Anschaulich kann man sagen: Eine kovalente Bindung entsteht nicht dadurch,
dass zwei Atome durch eine mechanische Stange verbunden wären. Sie entsteht
durch eine gemeinsame Elektronenverteilung, die die Atomkerne zusammenhält.

Ein einfaches Beispiel ist das Wasserstoffmolekül \( \mathrm{H_2} \). Zwei
Wasserstoffatome bringen jeweils ein Elektron mit. Nähern sich die Atome, so
überlagern sich ihre Atomorbitale. Es entsteht ein gemeinsames Molekülorbital,
in dem die Elektronendichte zwischen den beiden Kernen erhöht ist. Diese
erhöhte Elektronendichte stabilisiert das Molekül.

\begin{PhysicsBox}[Elektronendichte]
	\label{box:elektronendichte}
	Die Elektronendichte beschreibt die räumliche Verteilung der Elektronen in
	einem Atom, Molekül oder Festkörper.
	
	In Molekülen ist besonders wichtig, wo zwischen den Atomkernen erhöhte
	Elektronendichte auftritt. Dort kann eine bindende Wechselwirkung entstehen.
\end{PhysicsBox}

Nicht jede Überlagerung von Atomorbitalen führt zu einer stabilen Bindung.
Vereinfacht kann man zwischen bindenden und antibindenden Molekülorbitalen
unterscheiden. In einem bindenden Molekülorbital ist die Elektronendichte
zwischen den Kernen erhöht. Das stabilisiert das Molekül. In einem
antibindenden Molekülorbital ist die Elektronendichte zwischen den Kernen
vermindert. Das kann eine Bindung schwächen oder verhindern.

Für das Ziel dieses Buches ist vor allem der Grundgedanke wichtig: Die
chemische Bindung ist letztlich eine Folge der quantenmechanischen Struktur
der Elektronen. Elektronen bilden keine kleinen Verbindungshaken zwischen
Atomen. Sie erzeugen räumliche Verteilungen, die bestimmen, ob eine Verbindung
stabil ist.

Damit wird auch klar, warum die einfache Schalen- und Bindungsvorstellung nur
ein erster Schritt ist. Sie erklärt viele Grundtendenzen, aber die genauere
Beschreibung von Molekülen verlangt den Begriff der Orbitale und der
Elektronendichte.

Im nächsten Abschnitt wird dieser Gedanke auf Festkörper erweitert. Dort
gehören Elektronen nicht mehr nur zu einzelnen Atomen oder kleinen Molekülen,
sondern können sich über viele Atome hinweg ausbreiten. Daraus entstehen
Bänder, Leitfähigkeit und die elektronischen Eigenschaften von Metallen und
Halbleitern.